% =============================================================================
% Resumo Expandido --- Modalidade EXPERIÊNCIA EXITOSA
% Sessão de Temas Livres --- III Simpósio Internacional de Educação do Ensino
% Superior: inteligência artificial e ensino superior
% Publicação como suplemento na Revista CALIBRE.
% Formatação: Times New Roman, título 14pt, texto 12pt, espaçamento 1,5,
% máx. 1200 palavras, notas de rodapé 11pt, referências NBR 6023 (máx. 20).
% =============================================================================
\documentclass[12pt,a4paper]{article}

\usepackage[utf8]{inputenc}
\usepackage[T1]{fontenc}
\usepackage[brazil]{babel}
\usepackage{times}
\usepackage[top=3cm,left=3cm,bottom=2cm,right=2cm]{geometry}
\usepackage{setspace}
\onehalfspacing
\usepackage{indentfirst}

% Notas de rodapé em 11pt (conforme edital)
\renewcommand{\footnotesize}{\fontsize{11pt}{13pt}\selectfont}

% Bibliografia ABNT NBR 6023 --- hyperref ANTES de abntcite (ordem crítica)
\usepackage{hyperref}
\hypersetup{colorlinks=true,linkcolor=black,citecolor=black,urlcolor=blue}
\usepackage[alf]{abntcite}
\bibliographystyle{abnt-alf}
\providecommand{\abntrefinfo}[3]{}
\providecommand{\abntbstabout}[1]{}
\providecommand{\abntreprintinfo}[1]{\citeonline{#1}}

\setlength{\parindent}{1.25cm}
\setlength{\parskip}{0.2cm}

% Rótulos das seções da modalidade (texto em 12pt, em negrito)
\newcommand{\secao}[1]{\medskip\noindent\textbf{#1}\par\nobreak\vspace{-0.1cm}}

\pagestyle{empty}

\begin{document}

% ---------------------------------------------------------------------------
% TÍTULO (14pt) --- português e inglês, até 2 linhas cada
% ---------------------------------------------------------------------------
\begin{center}
{\fontsize{14pt}{17pt}\selectfont\bfseries
Inteligência Artificial no Apoio à Avaliação Formativa em Fábricas de
Software Acadêmicas: uma Experiência Exitosa\par}

\vspace{6pt}

{\fontsize{14pt}{17pt}\selectfont\bfseries\itshape
Artificial Intelligence Supporting Formative Assessment in Academic
Software Factories: a Successful Experience\par}

\vspace{12pt}

% ---------------------------------------------------------------------------
% AUTORES (instituição + titulação; e-mail do autor correspondente)
% ---------------------------------------------------------------------------
{\normalsize
Esp.~Danrley Pereira\footnote{Autor correspondente. E-mail:
\texttt{danrley.pereira@cs.udf.edu.br}} \quad
Dra.~Kerlla de Souza Luz\par}

\vspace{2pt}

{\small Centro Universitário do Distrito Federal (UDF) --- Curso de Ciência da
Computação, Brasília, Brasil\par}
\end{center}

\vspace{6pt}

\noindent\textbf{Palavras-chave:} Inteligência Artificial; Fábrica de Software
Acadêmica; Avaliação Formativa; \textit{Software Analytics}; Ensino Superior.

\vspace{4pt}

% ---------------------------------------------------------------------------
% SITUAÇÃO
% ---------------------------------------------------------------------------
\secao{Situação}

A Fábrica de Software Acadêmica LabTech, projeto de extensão do Centro
Universitário do Distrito Federal (UDF), insere estudantes dos cursos de
Tecnologia da Informação em projetos reais, com metodologias ágeis e
ferramentas da indústria, aproximando teoria e prática
\cite{cusumano1991japan,dos2021experiencia,oliveira2017conduccao}. O relato de
experiência do programa entre 2022 e 2024 \cite{pereira2026cadernos} evidenciou
ganhos consistentes de competências técnicas e interpessoais: em instrumentos
validados aplicados a quatorze participantes, as médias do
\textit{Entrepreneurship Competence Framework} (EntreComp)
\cite{bacigalupo2016entrecomp} situaram-se entre 3,51 e 3,79 (escala 1--4) e o
escore de autoeficácia geral \cite{sherer1982gse} concentrou-se na metade
superior da escala (média 79,5/100).

Contudo, o mesmo relato reafirmou uma limitação metodológica central: a
ausência de instrumentação contínua e de registros padronizados das atividades.
Sem visibilidade sistemática do andamento dos projetos, os docentes enfrentavam
assimetria de informação, dificuldade para avaliar objetivamente a colaboração
entre discentes e pouca capacidade de detectar precocemente estudantes ou
equipes em risco de baixo engajamento. A leitura manual de repositórios é
trabalhosa, e as estatísticas nativas do GitHub mostram-se insuficientes para a
decisão pedagógica \cite{hamer2021using}. Essa lacuna --- registrada por meio de
diários de bordo, entrevistas e análise de conteúdo --- motivou a adoção de
Inteligência Artificial (IA) e de \textit{analytics} automatizado como apoio à
avaliação formativa, dando origem à experiência aqui relatada.

% ---------------------------------------------------------------------------
% AÇÃO REALIZADA
% ---------------------------------------------------------------------------
\secao{Ação realizada}

Para endereçar essa necessidade, desenvolveu-se, no âmbito do Programa
Institucional de Bolsas de Iniciação Tecnológica e Inovação (PIBIT) do UDF, um
\textit{dashboard} de monitoramento de código aberto (licença GPL v3.0) para
acompanhamento contínuo de projetos em fábricas de software acadêmicas. A
plataforma extrai dados da API do GitHub --- \textit{commits}, \textit{pull
requests}, \textit{issues}, revisões e tráfego de repositórios --- e os organiza
em uma Arquitetura Medalhão de três camadas (\textit{Bronze}, dados brutos;
\textit{Silver}, dados normalizados; \textit{Gold}, métricas agregadas),
preservando a linhagem desde a origem até o indicador final e garantindo
reprodutibilidade.

Sobre essa base, calcula-se um conjunto de indicadores-chave de desempenho
(KPIs) técnicos (frequência de \textit{commits}, \textit{code churn},
\textit{lead time} de PR) e colaborativos (participação em revisões, rede de
colaboração entre autores e revisores, regularidade semanal e concentração de
contribuições). A seleção dos KPIs seguiu a abordagem \textit{Goal--Question--Metric}
\cite{basili1994gqm} e o processo de medição da norma ISO/IEC/IEEE 15939
\cite{iso15939}, com atenção explícita aos riscos de interpretação ingênua de
contagens brutas \cite{fenton2014metrics}.

O diferencial da experiência é o módulo de IA: a plataforma integra um modelo de
linguagem de grande escala (Google Gemini) \cite{googlegemini2024} que gera, em
linguagem natural, análises das contribuições de cada estudante, oferecendo ao
docente uma interpretação contextual dos dados quantitativos. O
\textit{dashboard} foi construído em ciclos iterativos de co-desenvolvimento
academia--academia \cite{marijan2020lessons} com uma docente parceira
(laboratório LabLivre/UnB), que definia, a cada iteração, quais dados
\textit{Silver} e \textit{Gold} eram necessários para acompanhar suas turmas; o
módulo de IA foi acrescentado justamente para atender à demanda por
interpretações contextuais. A solução opera com infraestrutura zero, via GitHub
Actions (\textit{pipeline} diário) e GitHub Pages, à semelhança de ferramentas
de monitoramento ágil \cite{ozkan2019agile,vega2022scrumwatch}.

% ---------------------------------------------------------------------------
% RESULTADOS OBTIDOS
% ---------------------------------------------------------------------------
\secao{Resultados obtidos}

O \textit{dashboard} foi implantado em duas instâncias independentes em
produção, em dois contextos pedagógicos reais: a disciplina de Métodos de
Desenvolvimento de Software (MDS) e o projeto de extensão LabTech, na UDF. Em
conjunto, as instâncias cobriram mais de uma centena de repositórios e cerca de
quatrocentos colaboradores externos acumulados ao longo dos ciclos da disciplina
de MDS --- essencialmente todos os discentes que por ela passaram.

A validação mais significativa foi de natureza pedagógica: ao final do semestre
2025/2, a docente parceira utilizou o \textit{dashboard}, apoiada pelas análises
geradas por IA, para avaliar seus alunos, analisando mais de setenta
repositórios de sua organização. Esse uso real confirmou que indicadores como o
ritmo inicial de contribuições (inclinação das primeiras oito semanas), a
regularidade semanal e a centralidade na rede de colaboração são suficientemente
acionáveis para a avaliação formativa.

Da percepção docente, registraram-se três efeitos principais. Primeiro, os
resumos em linguagem natural produzidos pela IA reduziram a leitura ingênua de
contagens brutas, fornecendo interpretação contextual das contribuições.
Segundo, os grafos de colaboração permitiram distinguir perfis de estudantes
\textit{integradores} (com muitos vínculos entre áreas) de \textit{especialistas}
(focados em uma trilha), enriquecendo a avaliação para além do volume de
\textit{commits}. Terceiro, a atualização semanal mostrou-se cadência suficiente
para o acompanhamento e a governança pedagógica. Esses achados convergem com os
ganhos de competência colaborativa já observados no relato de experiência
anterior (dimensão de criação de valor colaborativa do EntreComp~=~3,79), agora
instrumentados de forma contínua e automatizada.

% ---------------------------------------------------------------------------
% CONCLUSÃO
% ---------------------------------------------------------------------------
\secao{Conclusão}

A experiência demonstra que a Inteligência Artificial, combinada a um
\textit{pipeline} aberto de \textit{software analytics}, pode apoiar de forma
eficaz a avaliação formativa em fábricas de software acadêmicas, qualificando o
\textit{feedback} docente e democratizando o acesso a \textit{analytics} ---
uma vez que a solução é gratuita, de código aberto e opera sem infraestrutura
própria. A principal contribuição é mostrar que práticas de engenharia de dados
da indústria, aliadas a um modelo de linguagem, transformam dados dispersos de
repositórios em informação pedagógica acionável. Ressalva-se, porém, que a IA
atuou como apoio --- e não substituto --- da decisão docente: combinações
automáticas de KPIs não devem ditar a nota final. Como perspectivas futuras,
preveem-se a incorporação de predições baseadas em aprendizado de máquina para a
identificação precoce de equipes em risco, a expansão para novas instituições e
disciplinas e a integração com plataformas além do GitHub.

% ---------------------------------------------------------------------------
% REFERÊNCIAS (NBR 6023)
% ---------------------------------------------------------------------------
\begin{singlespace}
\small
\bibliography{resumo-simposio-ia}
\end{singlespace}

\end{document}
